\section{Introduction}

The goal of this thesis can be boild down to be just one proof. So this begs the question 
\textbf{ Why prove things, and what should a proof do?}


On November~20, 2026, Levent Alpöge released a proof of the Riemann
Hypothesis as a Lean library comprising 17~million lines of code
\cite{alpoege2026}. This raises the question of why we, as
mathematicians, want proofs at all. Confronted with the (then
hypothetical) existence of a proof of the Riemann Hypothesis that is
unreadable for humans, Alex Kontorovich said \cite[8:30\,min ff.]{3b1b2025}:

\begin{quote}
\enquote{The reason proof is [\dots] our epistemology, [\dots] [i.e.\,] the mechanism
by which we arrive at truth, is because we've been burned so long on
intuitive arguments. And you probably know a thousand different, you
know, fake proofs of $1=0$ and mistakes in Euclid's \emph{Elements}.
Human intuition is incredible and it's also very faulty. And for that
reason, we've decided that we need to have rigorous proof within an
axiomatic system as a mechanism for checking our understanding. But
that's not understanding itself. You know, you can understand a proof
in a picture. You can say, I see very clearly what's going on. Yes, I
have the technique to convert what I see clearly into an argument in
language, and that's like the gold standard. But the understanding
happens somewhere else.}
\end{quote}

In the same conversation, Ankur Moitra added \cite[10:59\,min ff.]{3b1b2025}:

\begin{quote}
\enquote{Historically, when something becomes cheap or commoditized in
society, the things we value become orthogonalized with respect to
those directions. [\dots] But the question is what is going to be
commoditized, right? What is going to become cheap and less valued than
it was before? And what are the orthogonal directions that are going to
survive for it, right? And it does seem to me that certain types of
problem solving are going to become cheaper. [\dots]}
\end{quote}

From this, my view is that the goal of mathematics is to expand human
understanding of logical and structural relationships. In the past, a
formal proof was the best way to bring understanding and correctness
together. Today, however, AI systems publish results that are almost
unreadable but formally verifiable; on October~26, 2026, for instance,
OpenAI published 721~manuscripts in a single day \cite{openai2026}.
The epistemology of mathematics therefore needs to rethink how
understanding and correctness can be brought together, in my opinion by
placing more value on intuitive explanations and investing more time
in them.